\documentclass[dvipdfm,cjk,14pt,hyperref={bookmarks=false,compress,slidestop}]{beamer}
\AtBeginDvi{\special{pdf:tounicode EUC-UCS2}}
\renewcommand{\kanjifamilydefault}{\mgdefault}
\renewcommand{\baselinestretch}{1.4}
\usefonttheme{professionalfonts}
\usetheme{GrayGradient}
\usepackage[T1]{fontenc}
\usepackage[utf8]{inputenc}
\usepackage{lmodern}
\usepackage[sc]{mathpazo}
\usepackage[scaled]{helvet}
\usepackage[scaled]{beramono}
\usepackage[deluxe,expert]{otf}
\usepackage{textcomp,okumacro}
\usepackage{url}
\graphicspath{{../images/}}
\author{金子達哉 \texttt{(id:catatsuy)}}
\title{Scala 初心者が Scala を勉強しようとして}
\date{}
\begin{document}

\begin{frame}
 \titlepage
\end{frame}

\begin{frame}
 \frametitle{自己紹介}
 \begin{itemize}
  \item 金子達哉
  \item はてな ID: \texttt{catatsuy}
  \item twitter: \texttt{catatsuy}
 \end{itemize}

\begin{picture}(0,0)(0,0)
 \put(200,0){\includegraphics[clip, height=35truemm]{catatsuy}}
\end{picture}
 \vspace{-5pt}
 URL:
 \begin{itemize}
  \item \url{http://www.catatsuy.org}
  \item \url{http://blog.catatsuy.org}
  \item \url{https://matw.co}
 \end{itemize}
\end{frame}

\begin{frame}
 \frametitle{所属}
 \begin{itemize}
  \item 東京工業大学
  \item 情報工学科 4 年（9 月卒業予定）
  \item 吉瀬研究室\\
        \begin{itemize}
         \item コンピュータアーキテクチャ
        \end{itemize}
 \end{itemize}
\end{frame}

\begin{frame}
 \frametitle{就職活動}
 \begin{itemize}
  \item はてなインターン 2012
  \item pixiv インターン
 \end{itemize}
 \vspace{-20pt}
 \begin{center}
  \begin{tabular}{cc}
   \includegraphics[clip, height=38truemm]{pixiv} & \includegraphics[clip, height=38truemm]{hatena} \\ 
  \end{tabular}
 \end{center}
\end{frame}

\begin{frame}[plain]
 \begin{center}
  \LARGE
  \includegraphics[clip, height=58truemm]{pixivall}\\
  \bfseries\gtfamily 10 月から \includegraphics[clip,
  height=25pt]{pixiv_logo} \hspace{2pt}  へ！
 \end{center}
\end{frame}

\begin{frame}[plain]
 \begin{center}
  \Huge
  \textbf{WARNING}
 \end{center}
\end{frame}

\begin{frame}[plain]
 \begin{center}
  \LARGE
  このプレゼンはゆるふわです\\
  コードは一行も出てきません\\
  リラックスして聞きましょう
 \end{center}
\end{frame}

\begin{frame}
 \frametitle{何でいるの？}
 \begin{picture}(0,0)(0,0)
  \put(0,5){\includegraphics[clip,
  height=30truemm]{tweet_288671004547502080.png}}
  \put(150,-35){\includegraphics[clip,
  height=30truemm]{tweet_296299794069061633.png}}
  \put(0,-95){\includegraphics[clip,
  height=30truemm]{tweet_296299960557793280.png}}
 \end{picture}
\end{frame}

\begin{frame}
 \frametitle{お手柔らかにお願いします}
 \begin{picture}(0,0)(0,0)
  \put(40,-60){\includegraphics[clip,
  width=90truemm]{tweet_296301154684514304.png}}
 \end{picture}
\end{frame}

\begin{frame}
 \frametitle{Scala に興味を持った理由}
 \begin{itemize}
  \item 関数型言語
  \item 型推論
  \item オブジェクト指向
  \item シンプルな記述
  \item JVM の上で動く
 \end{itemize}
\end{frame}

\begin{frame}[plain]
 \begin{itemize}
  \item 普段使っている言語を部分的に関数型言語みたいに使えないのか
  \item それができれば無理して使う必要はない？
 \end{itemize}
\end{frame}

\begin{frame}[containsverbatim]
 \frametitle{最高にやばい言語 PHP}
\small
 \begin{verbatim}
array array_map (callable $callback, array $arr1)
bool array_walk (array &$array, callable $funcname)
array array_filter (array $input, callable $callback)
 \end{verbatim}

\begin{center}
\Large
 嘘みたいだろ？ 同じ言語なんだぜ\\
これで…
\end{center}

\end{frame}

\begin{frame}
 \frametitle{PHP の問題点}
 \begin{itemize}
  \item PHP5.3 から無名関数が使えるように！
  \item 組込み関数の使い勝手最悪で使う気にならない
  \item \texttt{0 == 'hoge'} 問題が発生する
  \item 基本的な部分でも結構はまる
 \end{itemize}
\end{frame}

\begin{frame}
 \frametitle{コップ本を購入しました！}
 \begin{picture}(0,0)(0,0)
  \put(160,-75){\includegraphics[clip,
  width=55truemm]{scala_cop_book.jpg}}
 \end{picture}
 Scalaスケーラブル\\ プログラミング第2版
\end{frame}

\begin{frame}[plain]
 \begin{center}
  \LARGE よくわからない…
 \end{center}
\end{frame}

\begin{frame}
 \frametitle{分からない理由}
 \begin{itemize}
  \item 関数型・オブジェクト指向型・手続き型への理解が浅い
  \item 用語・メリットが分からない
  \item 関数型の勉強をしたいがこれでは勉強にならない
 \end{itemize}
\end{frame}

\begin{frame}
 \frametitle{すごい H 本}
 \begin{picture}(0,0)(0,0)
  \put(180,-75){\includegraphics[clip,
  width=55truemm]{sugoi_h.jpg}}
 \end{picture}
 \noindent 関数型言語の勉強のために\\ すごい H 本を購入
\end{frame}

\begin{frame}
 \frametitle{すごい H 本で}
 \begin{itemize}
  \item 関数型言語の用語や威力など勉強できる
  \item 関数型言語の経験が浅い人はコップ本読んでも意味ない
  \item まずはすごい H 本！
 \end{itemize}
\end{frame}

\begin{frame}
 \frametitle{まとめ}
 \begin{itemize}
  \item まずすごい H 本読みましょう
  \item 電子書籍あります
  \item そのあとコップ本読みましょう
  \item Scala 勉強します
 \end{itemize}
\end{frame}


\end{document}
